\section{The mapping heuristic}
\label{sec:method}

\begin{inbrief}
A design point runs seven steps (\cref{fig:pipeline}): partitioning into
sub-chains, allocation by bin packing or by a communication-aware rule,
core consolidation, priority assignment, relocation moves for
unschedulable cores, a deadline-driven local search, and the final
analysis. Contention is handled by a fixed point on the number of cores
of each ECU, and an outer loop can size the platform for a target
margin. Every step is deterministic.
\end{inbrief}

A \emph{design point} is one run of the method for one application and
one set of parameters (\cref{tab:params}). The decomposition follows the
objectives. Partitioning, allocation and consolidation decide the cores
(O1). Priority assignment and the moves of step 5 make every core
schedulable (C7); step 5 adds no core to the platform, but may put an
unused core back into use, so it can raise $M$. The local search of step 6
improves (O2) without changing $M$.

\begin{figure*}[t]
\centering
\resizebox{\textwidth}{!}{%
\begin{tikzpicture}[node distance=3.0mm and 3.0mm,
	st/.style={stage, minimum width=2.05cm, minimum height=0.95cm}]
\node[st, fill=cpcu!60] (app) {application\\[-1pt]{\scriptsize task-chains $\Chains$}};
\node[st, right=of app] (part) {\textbf{1} partitioning\\[-1pt]{\scriptsize sub-chains, \cref{alg:partition}}};
\node[st, right=of part] (assign) {\textbf{2} allocation\\[-1pt]{\scriptsize bin packing / \CADU}};
\node[st, right=of assign] (cons) {\textbf{3} consolidation\\[-1pt]{\scriptsize \cref{alg:consolidate}}};
\node[st, right=of cons] (sched) {\textbf{4} priorities\\[-1pt]{\scriptsize \cref{alg:audsley}}};
\node[st, right=of sched] (place) {\textbf{5} relocation\\[-1pt]moves};
\node[st, right=of place] (repair) {\textbf{6} deadline\\[-1pt]local search};
\node[st, right=of repair, fill=czcu!70] (out) {\textbf{7} analysis\\[-1pt]{\scriptsize metrics, $M_{\mathrm{LB}}$}};
\foreach \a/\b in {app/part,part/assign,assign/cons,cons/sched,sched/place,place/repair,repair/out}
	\draw[arrow] (\a) -- (\b);
\node[note, below=1.2mm of cons] {decide $M$ (O1)};
\node[note, below=1.2mm of place] {schedulable (C7)};
\node[note, below=1.2mm of repair] {improves $N$ (O2)};
\draw[arrow, dashed, color=accent, rounded corners=3pt]
	(out.north) -- ++(0,0.42) -| node[note, pos=0.25, above, color=accent]
	{sizing for $\lambda^\star$: bisection on $\varphi$ (\cref{sec:sizing})}
	(part.north);
\draw[arrow, dashed, color=cblue, rounded corners=3pt]
	(assign.south) -- ++(0,-0.62) -| node[note, pos=0.25, below, color=cblue]
	{contention: more cores on an ECU than assumed, run again (\cref{sec:contention-method})}
	(part.south);
\end{tikzpicture}}
\caption{The steps of a design point. The network analysis
(\cref{alg:network}) runs inside steps 6 and 7. The two loops are
optional: sizing for a margin, and the fixed point of the contention.}
\label{fig:pipeline}
\end{figure*}

\begin{table}[t]
\centering\footnotesize
\caption{Parameters of a design point and their defaults.}
\label{tab:params}
\setlength{\tabcolsep}{3pt}
\rowcolors{2}{brandlight}{white}
\begin{tabularx}{\columnwidth}{@{}llX@{}}
\toprule
\rowcolor{white}
parameter & default & meaning \\
\midrule
rule & \FFDU & allocation rule: \BFDU, \WFDU, \FFDU, \CADU \\
$\sigma$ & 0.4 & utilization limit of a sub-chain, on the core type it will get \\
$\varphi$ & 0.8 & utilization cap of a core \\
central cores & 8 big (100\,\%) & + little cores (40\,\%), none by default \\
zone cores & 4 MCU (50\,\%) & per zone; + lock-step cores, none \\
$c_{\LET}$ & \SI{50}{\micro s} & LET copy on the reference core \\
$M_k$ & unlimited & memory of each ECU \\
replicas & ECUs & separation policy (or cores) \\
affinity & on & try the cores of the task-chain first \\
consolidation & on & step 3 (not with \WFDU) \\
relocation & on & step 5 \\
swap moves & off & step 5, before relocation (\cref{sec:placement}) \\
local search & on & step 6 \\
$\ell$, $\eta$ & 0, 1000 & memory latency (ns); requests per ms of WCET \\
network & star, 1\,Gbit/s & topology, $R$; $R_b=R$ \\
shaping & CBS, 1 class & idle slope 75\,\%, no preemption; or TAS, ATS \\
FRER & off & ring or custom topology: replicated task-chains \\
margin & 0\,\% & WCETs increased before the mapping \\
$\lambda^\star$ & none & sizing target (\cref{sec:sizing}) \\
\bottomrule
\end{tabularx}
\end{table}

\paragraph{Preparation.} With a WCET margin of $m$\,\%, every reference
WCET first becomes $\min(\lceil C^{\mathrm{ref}}(100+m)/100\rceil,\,T-1)$,
and the platform is sized for the larger WCETs. \Cref{sec:sizing} offers
an alternative that sizes for the margin directly. With contention, the
tasks without a request count get $a_\tau=\lceil\eta
C^{\mathrm{ref}}_\tau/1000\rceil$ before the margin, which inflates
execution, not memory traffic.

\subsection{Step 1: partitioning}
\label{sec:partitioning}

Partitioning cuts each task-chain into sub-chains, the units of
allocation. Fewer sub-chains mean fewer communications; smaller sub-chains
pack better. A sub-chain must not be larger than what a core can take next
to others, so its utilization on the core type it will run on is limited
to $\sigma$ (minus 1\,\% kept for the LET task): on a core of speed $s$,
that is a utilization of $(\sigma-0.01)\,s/100$ of the reference core, as
\cref{alg:partition} computes. Its tasks must share their binding and
their allowed sets: a sensor driver never drags the computation into its
zone. Tasks that must be co-located (C2) are never separated, since they
have the same constraints (\cref{sec:problem}). \Cref{alg:partition} gives the procedure. A sub-chain of free
tasks is sized for the core it is most likely to get: the first host
(\cref{sec:other-architectures}) it may run on, else the first other ECU
it may run on, with the first type of that ECU it allows, in the order
big, little, MCU, lock-step. The \emph{fragmentation rate} is the
number of cuts over the number of possible cuts (one per pair of
consecutive tasks).

\begin{algorithm2e}[t]
\caption{\textsf{Partition}$(\xi,\sigma)$}
\label{alg:partition}
\Input{task-chain $\xi=(\tau_1..\tau_n)$, limit $\sigma$}
\Output{sub-chains of $\xi$, in chain order}
$F\gets\langle\rangle$ \tcp*{the open sub-chain}
\For{$i\gets1$ \KwTo $n$}{
	\If{$F\ne\langle\rangle$ \textbf{and} $(A_{\tau_i},\Theta_{\tau_i},\text{binding})\ne$ those of $F$}{
		close $F$ \tcp*{forced cut}}
	$g\gets F\ne\langle\rangle\wedge g_{\mathrm{last}(F)}$ \tcp*{glued to $F$}
	$\bar u\gets(\sigma-0.01)\,s_{\hat\theta}/100$, $\hat\theta$ the type $F$ will get\;
	\If{$U^{\mathrm{ref}}_{\tau_i}>\bar u$ \textbf{and not} $g$}{
		close $F$; $F\gets\langle\tau_i\rangle$\;
		\lIf{\textbf{not} $g_{\tau_i}$}{close $F$}
		\textbf{continue}\;}
	append $\tau_i$ to $F$\;
	\lIf{$g$}{\textbf{continue}}
	\If{$|F|>1$ \textbf{and} ($U^{\mathrm{ref}}(F)>\bar u$ \textbf{or} $F$ alone on a core is unschedulable)}{
		remove $\tau_i$; close $F$; $i\gets i-1$\;}
}
close $F$\;
\end{algorithm2e}

The schedulability test ``alone on a core'' in \cref{alg:partition} uses
chain order inside a pattern and rate-monotonic order between patterns,
with \cref{eq:rta}. The sub-chains of all task-chains are finally sorted by
decreasing reference utilization (stable sort), the order in which they are
allocated.

\subsection{Step 2: allocation}
\label{sec:assignment}

Sub-chains are placed one by one in that order (\cref{alg:assign}). The
load a sub-chain adds to a core is its utilization at the speed and type of
the core (\cref{eq:wcet}), plus the change of the LET load: $T^{\LET}$
becomes the GCD of the old and new periods (\cref{eq:let}). A core
\emph{can take} $f$ if (C1), (C4) and (C5) hold and the new load is within
$\varphi$ (C3).

\begin{algorithm2e}[t]
\caption{\textsf{Allocate}$(\Frags)$}
\label{alg:assign}
\Input{sub-chains by decreasing utilization; initial cores $\Cores_0$}
\Output{$\pi$, possibly with added cores, or \emph{infeasible}}
$\Cores\gets\Cores_0$; $\mathcal{R}_k\gets\sum_{f\in\Frags_k}\mu_f$ for all $k$ \tcp*{reserved memory}
\ForEach{$f\in\Frags$}{
	\For{$\textit{attempt}\gets1$ \KwTo $2$}{
		$p\gets\bot$\;
		\lIf{affinity}{$p\gets$ \textsf{Choose}$(f,\textsc{siblings})$}
		\lIf{$p=\bot$}{$p\gets$ \textsf{Choose}$(f,\textsc{preferred})$}
		\lIf{$p=\bot$}{$p\gets$ \textsf{Choose}$(f,\textsc{any})$}
		\lIf{$p\ne\bot$}{$\pi(f)\gets p$; \textbf{break}}
		\lIf{$\textit{attempt}=2$ \textbf{or} fixed hardware}{\Return \emph{infeasible}}
		$k\gets$ first ECU (hosts first) in $A_f$ whose memory, replicas and $M^{\max}_k$ admit $f$\;
		\lIf{none}{\Return \emph{infeasible}}
		add a core to $k$ of the added type of $k$, or, if $f$ may not use it, of the type in $\Types(k)\cap\Theta_f$ on which $f$ is lightest\;
	}
	\lIf{$f$ pinned to $k$}{$\mathcal{R}_k\gets\mathcal{R}_k-\mu_f$}
}
\end{algorithm2e}

\paragraph{Tiers.} \textsf{Choose}$(f,\textit{tier})$ considers the cores
that can take $f$, restricted by the tier:
\begin{itemize}
\item \textsc{siblings}: cores that already run another sub-chain of the
same task-chain (affinity; it saves a communication);
\item \textsc{preferred}: for a free sub-chain that may run on the central
computer, its cores only (except with \CADU); otherwise all allowed cores;
\item \textsc{any}: all allowed cores.
\end{itemize}
In every tier, for a free sub-chain that may run on the central computer,
the zone cores \emph{added} during the allocation are excluded: they were
added for the pinned sub-chains of their zone, and filling them with free
work would force yet more zone cores. Here and below, ``the central
computer'' stands for the hosts of \cref{sec:other-architectures} (the
central computers, or, on a platform without one, the ECUs with the
fastest cores), and ``zone cores'' for the cores of the other ECUs.

\paragraph{Rules.} Among the candidates, the rule picks:
\begin{itemize}
\item \BFDU\ (best fit decreasing utilization): the core left with the
least room;
\item \WFDU\ (worst fit): the core left with the most room;
\item \FFDU\ (first fit): the first core in creation order;
\item \CADU\ (communication-aware, new): the core that is lexicographically
best for (i) the number of inter-ECU communications the placement creates
with the neighbouring sub-chains of the task-chain, (ii) the number of
communications between cores of one ECU it creates, (iii) being a core of
the central computer, and (iv) the least room left.
\end{itemize}
For \CADU, the neighbours of $f$ are the previous and next sub-chains of
its task-chain. A neighbour already placed on core $q$ creates no
communication if $q=p$, a communication inside the ECU if $\kappa(q)=
\kappa(p)$, and an inter-ECU one otherwise. A neighbour not yet placed but
pinned to ECU $k'\ne\kappa(p)$ creates an inter-ECU communication. Bin
packing alone sends a free sub-chain to the central computer whenever a
core there can take it, unless affinity first finds a core that runs its
own task-chain. \CADU\ instead keeps a computation next to the sensor or actuator it
serves, when the zone controller has room for it.

\paragraph{Memory reservation.} A zone controller with a memory limit could
see its memory taken by free sub-chains placed early, and then fail to host
its own sensor and actuator drivers, which can run nowhere else. The
allocation therefore reserves, on every ECU $k$, the memory
$\mathcal{R}_k$ of the pinned sub-chains not yet placed. A free sub-chain
may use only $M_k-\mathcal{R}_k$, minus what is already used.

\begin{proposition}[Memory]
\label{prop:memory}
If $\sum_{f\in\Frags_k}\mu_f\le M_k$ for every ECU $k$, no pinned
sub-chain is ever refused by (C4) during \cref{alg:assign}.
\end{proposition}
\begin{proof}
Let $u_k$ be the memory used on $k$. The invariant $u_k+\mathcal{R}_k\le
M_k$ holds initially. Placing a pinned sub-chain on $k$ moves $\mu_f$ from
$\mathcal{R}_k$ to $u_k$; a free sub-chain is admitted only if $u_k+\mu_f+
\mathcal{R}_k\le M_k$. Both keep the invariant. When a pinned sub-chain $f$
of $k$ is placed, $\mu_f\le\mathcal{R}_k$, so $u_k+\mu_f\le M_k$.
\end{proof}

\subsection{Step 3: core consolidation}
\label{sec:consolidation}

Decreasing-utilization bin packing places each sub-chain once and never
revisits it. A core opened early for a large sub-chain can stay lightly
loaded, while the smaller sub-chains allocated after it fill the gaps that
the pinned work leaves on other cores. The exact formulation of
\cref{sec:exact} showed that this costs one to three cores on small and
constrained instances. Consolidation (\cref{alg:consolidate}) removes this
loss. It takes the cores in use from the lightest, and tries to move all
their sub-chains to other cores in use. A move must respect (C1)--(C5) and
the cap $\varphi$ (and, as in the allocation, a free sub-chain does not go
to a zone core added during the allocation), and the target core must
stay schedulable with the
first-fit pass of \cref{alg:audsley}. Among the targets, the sub-chain goes
to the one that leaves the fewest inter-ECU communications with its
neighbours in the task-chain, then to the fullest, so that consolidation
does not undo the work of affinity and of \CADU. The moves are kept only if
the whole core is emptied; otherwise the core is left as it was.
Consolidation is not applied to \WFDU, which spreads the load on purpose.

\begin{algorithm2e}[t]
\caption{\textsf{Consolidate}$(\pi)$}
\label{alg:consolidate}
\Repeat{no core was emptied}{
	\ForEach{core $p$ in use, by increasing load}{
		$\pi'\gets\pi$\;
		\ForEach{sub-chain $f$ of $p$, by decreasing utilization}{
			$Q\gets$ cores $q\ne p$ in use that can take $f$ in $\pi'$\;
			sort $Q$ by (inter-ECU communications of $f$ on $q$, room left)\;
			$q\gets$ first $q\in Q$ schedulable with $f$ (first-fit priorities)\;
			\lIf{none}{\textbf{next} $p$ \tcp*[f]{keep $\pi$}}
			move $f$ from $p$ to $q$ in $\pi'$\;
		}
		$\pi\gets\pi'$ \tcp*{$p$ is empty}
		\textbf{restart}\;
	}
}
\end{algorithm2e}

Each kept round empties one core, so there are at most $M$ rounds, and
every move keeps (C1)--(C5) and (C3). Schedulability is checked on every
target with the optimal test of \cref{thm:audsley}, so consolidation never
makes a schedulable core unschedulable.

\subsection{Contention: the number of competing cores}
\label{sec:contention-method}

With contention, the WCETs depend on the number $n_k$ of cores of ECU $k$
that compete for its memory, and $n_k$ depends on the allocation. The
method resolves this circularity as follows.
\begin{enumerate}
\item The allocation starts with $n_k$ equal to the number of cores of $k$
in the architecture. A core added to $k$ during the allocation gets the
current count.
\item If the allocation took more cores on some ECU than assumed, the
design point is run again from step 1 with that many. More cores mean more
contention, longer WCETs and possibly more cores, so this is a fixed
point. It stops when the counts no longer grow, after at most eight runs;
otherwise the design point is declared infeasible.
\item After step 6, only the cores that run tasks compete: $n_k$ becomes
the number of cores of $k$ in use. The WCETs can only decrease, and the
response times are computed again with the same priorities.
\end{enumerate}

\begin{proposition}[Contention]
\label{prop:contention}
In the final mapping, every WCET is at least the one of \cref{eq:wcet}
with $n_k$ the number of cores of $k$ that run tasks, and every core
declared schedulable is schedulable with those WCETs.
\end{proposition}
\begin{proof}
At the end of the fixed point, the count assumed for every ECU is at
least the number of its cores, so every WCET used by steps 2 to 6 is at
least the final one. Item 3 of the list above recomputes the WCETs with the
number of cores in use, which no later step changes, and recomputes the
response times with the final priorities. The response time of
\cref{eq:rta} is monotone in the WCETs.
\end{proof}

\subsection{Step 4: priorities}

Every core with tasks gets its priorities from \cref{alg:audsley} with the
robust choice. Bin packing only checks loads, so a core within $\varphi$
can still be unschedulable: the automotive periods are not harmonic
(2\,ms and 5\,ms do not divide each other). Such cores are repaired in
step 5.

\subsection{Step 5: relocation moves}
\label{sec:placement}

This step moves sub-chains away from unschedulable cores, without adding
cores to the platform. The sub-chains of the unschedulable cores are
visited one at a time, the cores in their order and the sub-chains in the
order of their core. A sub-chain $f$ whose core has become schedulable is
skipped. Otherwise $f$ moves to the core that (i) satisfies (C1), (C4),
(C5) and the cap $\varphi$ with it, (ii) is schedulable, (iii) remains
schedulable with $f$ after \cref{alg:audsley}, and (iv) keeps the most
room. Cores without tasks (initial cores not used yet, or cores emptied
by step 3) are tried only when no loaded core qualifies, which is how
this step can raise $M$. The source core then gets new priorities.

\begin{remark}[Swap moves]
An earlier version of this step first exchanged pairs of sub-chains
between two unschedulable cores, and kept an exchange that made one of
them schedulable. The tool still offers these swap moves as an option, off
by default. Of the 9600 design points of \cref{sec:evaluation}, 9516
have an admissible mapping, and 84 of these have unschedulable cores
after step 4. The relocation moves, where enabled, made all of them
schedulable. In the variant of the ablation study with swap moves
(\cref{sec:rq8}), the swap moves made none of these cores schedulable.
\end{remark}

\subsection{Step 6: deadline-driven local search}
\label{sec:repair}

A task-chain whose tasks are all schedulable can still miss its
end-to-end deadline, because every communication costs LET sampling delays
and, between ECUs, a network transfer. A task-chain can also cross an
overloaded port. The local search removes communications of such
task-chains (\cref{alg:repair}). Around each communication it tries two
moves (\cref{fig:repair}): \emph{pull back} moves the first sub-chain after
the communication onto the core before it, and \emph{push on} moves the
last sub-chain before the communication onto the core after it. A move
must respect (C1)--(C5), leave the target core schedulable, and leave the
source core schedulable (or no worse). The source keeps its priorities,
because removing tasks cannot increase a response time. A move is
\emph{kept} only if the number of task-chains that meet their deadline,
over the whole application, strictly increases. The count must be global:
a move changes the response times of two cores and the loads of the ports,
and hence other task-chains. Trials are screened with first-fit priorities
on the target; a promising move gets the robust priorities and is counted
again.

\begin{algorithm2e}[t]
\caption{\textsf{LocalSearch}$(\pi)$}
\label{alg:repair}
$\mathcal{S}\gets$ \textsf{AnalysisInit}$(\pi)$; $N\gets N(\mathcal{S})$\;
\For{$\textit{pass}\gets1$ \KwTo $8$}{
	$\textit{kept}\gets$ false\;
	\ForEach{task-chain $\xi$ schedulable but late in $\mathcal{S}$}{
		\ForEach{communication $h$ of $\xi$, move $\mu\in\{$pull back, push on$\}$}{
			\lIf{$\mu$ violates (C1)--(C5), or the target is unschedulable with first-fit priorities}{\textbf{continue}}
			apply $\mu$ (cores $p,q$ changed); $n\gets$ \textsf{Update}$(\mathcal{S},p,q)$\;
			\lIf{$n>N$}{robust priorities on the target; $n\gets$ \textsf{Update}$(\mathcal{S},q)$}
			\eIf{$n>N$}{$N\gets n$; $\textit{kept}\gets$ true; restart with the same $\xi$\;}
			{undo $\mu$; \textsf{Update}$(\mathcal{S},p,q)$\;}
		}
	}
	\lIf{\textbf{not} \textit{kept}}{\textbf{break}}
}
\end{algorithm2e}

\begin{figure}[t]
\centering
\begin{tikzpicture}[node distance=1.4mm]
\node[task, fill=czcu] (a) {$f_0$};
\node[draw=cmuted, dashed, rounded corners=2pt, fit=(a), label={[note]above:ZCU2}] (c1) {};
\node[task, right=14mm of a] (b) {$f_1$};
\node[draw=cmuted, dashed, rounded corners=2pt, fit=(b), label={[note]above:central, core 0}] (c2) {};
\node[task, right=20mm of b] (c) {$f_2$};
\node[draw=cmuted, dashed, rounded corners=2pt, fit=(c), label={[note]above:central, core 3}] (c3) {};
\draw[arrow, color=corange] (c1) -- (c2);
\draw[arrow, color=corange] (c2) -- node[note, below] {communication} (c3);
\draw[arrow, color=cblue, out=125, in=55, looseness=1.1] ([yshift=5mm]c.north) to node[note, above] {pull back $f_2$} ([yshift=5mm]b.north);
\draw[arrow, color=cgreen, out=-55, in=-125, looseness=1.1] ([yshift=-1mm]b.south) to node[note, below] {push on $f_1$} ([yshift=-1mm]c.south);
\end{tikzpicture}
\caption{The two moves tried around a communication. Around the
communication from ZCU2, only the pull back of $f_1$ exists: $f_0$ is
pinned to its zone.}
\label{fig:repair}
\end{figure}

\paragraph{Incremental analysis.} Counting the task-chains that meet their
deadline from scratch after every trial costs the whole communication
analysis: collecting and sorting all sub-chains, rebuilding all segments
and analysing the network. The local search instead keeps a state
$\mathcal{S}$: the descriptor of every sub-chain (its core, ECU, the
response times that \cref{eq:delay,eq:reaction} need, and the LET
parameters of its core), the segments of every task-chain, and the
summary. \textsf{Update}$(\mathcal{S},p,q)$ recomputes the descriptors of
the sub-chains now on $p$ and $q$, and rebuilds only the segments of their
task-chains. It then reruns \cref{alg:network} on all flows, whose cost is
linear in the number of flows, and sums the bounds again.

\begin{proposition}[Exactness]
\label{prop:incremental}
After any sequence of moves, each followed by \textsf{Update} on the cores
it changed, $\mathcal{S}$ equals the result of a full analysis of the
current mapping.
\end{proposition}
\begin{proof}
A sub-chain descriptor depends only on the core that holds the sub-chain
(the response times of its tasks and the LET task of the core). A move
changes only the two cores it touches, so the descriptors of the
sub-chains on all other cores are unchanged, and the sub-chains that
changed core are now on $p$ or $q$. The segments of a task-chain depend only
on the descriptors of its sub-chains; the network analysis and the sums are
recomputed from all segments. The full analysis is a deterministic function
of the same data, in the same order.
\end{proof}
Undoing a rejected move is a move too, and is followed by \textsf{Update}.
The regression tests check \cref{prop:incremental} after random moves.

\subsection{Step 7: analysis and metrics}

Finally, the segments of all task-chains are built, the flows are analysed
together (\cref{alg:network}), and the bounds of \cref{prop:e2e} are
computed. The mapping is verified: no core above $\varphi$, and priorities
in precedence order within every pattern. The method also computes the
critical scaling factor of every core
(\cref{sec:sensitivity}) and the lower bound $M_{\mathrm{LB}}$
(\cref{prop:lb}).

\begin{keybox}{What a design point guarantees}
\textbf{Guaranteed}, under the model of \cref{sec:model}: every mapping
returned meets (C1)--(C6); every core either meets (C7) or is reported as
unschedulable; every response time (\cref{eq:rta}) reported within its
period, every port delay (\cref{sec:nc}), and every end-to-end delay and
maximum reaction time (\cref{prop:e2e}) reported is an upper bound; a
task above its period is reported without a response time, and a port
loaded above its reserved rate, or crossed by a flow of such a task, is
reported as giving no bound, and so is every task-chain that depends on
them; $M_{\mathrm{LB}}$ is a lower
bound on the cores of any mapping of the same sub-chains
(\cref{prop:lb}).
\textbf{Not guaranteed}: that $M$ is optimal (\cref{sec:exact} measures
the gap); that every deadline is met, since deadlines are an objective
(O2) and the late task-chains are reported; and anything outside the
model: WCETs, request counts and message sizes are inputs, and the
bounds across ECUs assume clocks that run at the same rate, though not
with aligned phases (\cref{sec:discussion}).
\end{keybox}

\subsection{Sizing for a margin}
\label{sec:sizing}

A WCET margin at design time inflates every WCET and hopes that the
resulting mapping keeps enough slack. The sizing variant of the problem
asks for $\lambda(\pi)\ge\lambda^\star$ directly. A lower utilization cap
leaves more room on every core and, in general, a larger scaling factor.
The method therefore searches for the largest $\varphi$ that reaches the
target. It runs the design point at the given $\varphi_0$ and stops if
$\lambda\ge\lambda^\star$. Otherwise it runs at $\varphi=\sigma$, the
smallest value at which sub-chains still fit; if that run misses the target
too, it is returned with the flag \emph{target missed}. Otherwise, seven
bisection steps on $[\sigma,\varphi_0)$ keep the largest $\varphi$ that
reached the target, so at most nine design points are run. The scaling
factor is not exactly monotone in $\varphi$, because the mapping changes
with it, so the result is the best $\varphi$ found, not a guaranteed
maximum. A run reaches the target when its scaling factor is at least
$\lambda^\star(1-10^{-3})$ (a tolerance for floating point) and every core
is schedulable.

\subsection{Complexity}

Let $n$ be the number of tasks, $N$ of sub-chains and $P$ of cores.
Partitioning is linear in $n$ times the cost of a response-time analysis
of one sub-chain. Allocation evaluates at most $P$ cores per sub-chain and
tier, $O(NP)$. A round of consolidation tries at most $P$ cores, each
costing a copy of the cores and one priority assignment per moved
sub-chain and target. Relocation tries at most $P$ targets for each
sub-chain of an unschedulable core, a set that is empty in most runs at
$\varphi\le0.8$. The
local search tries $O(H)$ moves per pass for $H$ communications of late
task-chains, each costing two priority assignments and one incremental
update, linear in the flows. \Cref{sec:scalability} reports the measured
growth.
